\section{Dataset}
\label{sec:dataset}

SnatchPhaseBench comprises annotated monocular snatch attempts with per-frame phase labels, segment interval annotations, and precomputed MediaPipe Pose Landmarker keypoints.
This section summarizes verified dataset statistics and preprocessing outputs; items requiring legal confirmation or additional metadata are marked explicitly.

\subsection{Dataset overview}
\label{sec:dataset:overview}

The dataset contains \numvideos{} snatch attempts performed by \numathletes{} athletes.
Each attempt is represented as a video file path key (raw videos are not bundled in the audited export), frame-wise phase annotations, optional segment intervals, and a CSV of pose keypoints aligned by frame index.
Seven phase classes are used for supervised learning; an additional \texttt{unlabeled} class appears in annotation files but is excluded from training windows when the center frame is unlabeled (\Cref{sec:dataset:preprocessing}).

After sliding-window construction, the processed dataset contains \numwindows{} samples of shape $(\windowsize, \featuredim)$ stored as \texttt{X.npy} with integer labels \texttt{y.npy}.
Independent rebuilding from annotation and keypoint CSV files reproduced the SHA-256 checksums recorded in the source manifest for \texttt{X.npy} and \texttt{y.npy}.

\subsection{Video collection and clip definition}
\label{sec:dataset:videos}

Each record corresponds to a single snatch attempt clip referenced by a relative path \texttt{video\_relpath} (e.g.\ \texttt{athlete\_folder/attempt.mp4}).
The audited repository includes 208 unique \texttt{video\_relpath} entries consistent with 208 keypoint CSV files.

\textsc{TODO (evidence required):} Document competition or training source, recording dates, camera viewpoints, native frame rate, resolution, and barbell visibility criteria.
\textsc{TODO (legal):} Confirm redistribution rights before public release.

\subsection{Phase taxonomy}
\label{sec:dataset:taxonomy}

Phase labels follow a seven-class snatch decomposition used in the source thesis pipeline.
\Cref{tab:phase_taxonomy} lists class identifiers; operational biomechanical definitions (joint-angle or barbell-height rules) were not independently verified during reproduction.

\begin{table}[t]
  \centering
  \caption{Phase taxonomy for supervised learning. Phase~0 is excluded when it appears at the center of a training window.}
  \label{tab:phase_taxonomy}
  \begin{tabular}{clp{0.48\linewidth}}
    \toprule
    ID & Phase name & Notes \\
    \midrule
    0 & \texttt{unlabeled} & Annotation padding / undefined frames; dropped at window center \\
    1 & \texttt{setup} & Pre-lift positioning \\
    2 & \texttt{first\_pull} & Initial barbell separation from platform \\
    3 & \texttt{transition} & Knee--hip coordination change between pulls \\
    4 & \texttt{second\_pull} & Explosive extension phase \\
    5 & \texttt{turnover} & Barbell ascent and arm turnover \\
    6 & \texttt{catch} & Receipt of the barbell overhead \\
    7 & \texttt{recovery} & Stabilization and stand-up \\
    \bottomrule
  \end{tabular}
\end{table}

\textsc{TODO:} Insert biomechanical definitions with citations once confirmed with domain experts; add figure illustrating phases on exemplar lift.

\subsection{Annotation strategy}
\label{sec:dataset:annotations}

Two annotation granularities are provided:
\begin{itemize}
  \item \textbf{Frame labels} (\texttt{master\_frame\_labels.csv}): one row per $(\text{video}, \text{frame})$ with \texttt{phase\_id} and \texttt{phase\_name}.
        The audited export contains 35{,}825 frame rows with no duplicate $(\text{video}, \text{frame})$ keys.
  \item \textbf{Segment labels} (\texttt{master\_segment\_labels.csv} and per-video CSV files): contiguous intervals per phase suitable for segment-level evaluation.
        856 per-video segment files are present in the audited export.
\end{itemize}

Annotation was performed using an interactive tool in the source project (\texttt{annotate\_phases.py} in the student repository).
\textsc{TODO:} Document annotator background, instructions, adjudication protocol, and inter-annotator agreement (not measured in current artifact).

\subsection{Pose representation}
\label{sec:dataset:pose}

Pose is represented using MediaPipe Pose Landmarker with \numlandmarks{} body landmarks.
Each frame stores normalized image coordinates $(x, y, z)$, yielding $\featuredim = 3 \times \numlandmarks$ numeric features per frame.
Visibility scores are available in keypoint CSV headers but are \emph{not} used by the frozen baseline pipeline (\texttt{use\_visibility: false}).

Missing values are handled via interpolation followed by forward and backward filling along each coordinate stream (verified in the frozen preprocessing implementation).
\textsc{TODO:} Report MediaPipe model version (\texttt{pose\_landmarker\_full.task}) and extraction settings once the binary task file is available locally (currently an incomplete export pointer in one snapshot).

\subsection{Train, validation, and test splits}
\label{sec:dataset:splits}

Generalization is evaluated under a fixed \emph{athlete-level} split: all attempts from a given athlete belong exclusively to training, validation, or test.
This prevents near-duplicate lifting styles from leaking across splits through highly correlated windows (stride~1; \Cref{sec:protocol:overlap}).

\Cref{tab:split_stats} summarizes verified split counts from \texttt{athlete\_split.json}.
No athlete or video identifier appears in more than one split (automated validation test in the canonical repository).

\begin{table}[t]
  \centering
  \caption{Verified athlete-level data split. Window counts derive from rebuilt \texttt{meta.csv}.}
  \label{tab:split_stats}
  \begin{tabular}{lrrr}
    \toprule
    Split & Athletes & Videos & Windows \\
    \midrule
    Train & 49 & 145 & 14{,}140 \\
    Validation & 10 & 30 & 3{,}232 \\
    Test & 11 & 33 & 3{,}877 \\
    \midrule
    Total & 70 & 208 & 21{,}249 \\
    \bottomrule
  \end{tabular}
\end{table}

\textsc{TODO:} Report athlete demographics (sex, weight class, performance level) if available and cleared for publication.
\textsc{TODO:} Evaluate multi-fold cross-validation as robustness analysis (planned; not yet executed).

\subsection{Preprocessing and generated tensors}
\label{sec:dataset:preprocessing}

The frozen preprocessing pipeline (\Cref{sec:methods:preprocessing}) constructs sliding windows of length $\windowsize = 31$ with stride~1.
Each window receives the phase label at its center frame; windows whose center is \texttt{unlabeled} are discarded.

Generated artifacts per rebuild:
\begin{itemize}
  \item \texttt{X.npy} $\in \mathbb{R}^{\numwindows \times \windowsize \times \featuredim}$ (\texttt{float32})
  \item \texttt{y.npy} $\in \mathbb{Z}^{\numwindows}$ (\texttt{int64})
  \item \texttt{meta.csv} with columns linking each window to \texttt{video\_relpath}, frame indices, athlete split, and center \texttt{phase\_id}
  \item \texttt{label\_map.csv} mapping phase IDs to names (eight entries including \texttt{unlabeled})
\end{itemize}

\Cref{tab:class_distribution} reports the verified window counts per phase in the rebuilt dataset.
Transition and \texttt{second\_pull} remain the least frequent actionable phases, suggesting class imbalance relevant to metric interpretation.

\begin{table}[t]
  \centering
  \caption{Verified window counts per phase after center-frame labeling (rebuilt dataset).}
  \label{tab:class_distribution}
  \begin{tabular}{lc}
    \toprule
    Phase & Windows \\
    \midrule
    setup & 6{,}017 \\
    first\_pull & 2{,}434 \\
    transition & 673 \\
    second\_pull & 1{,}245 \\
    turnover & 1{,}903 \\
    catch & 1{,}912 \\
    recovery & 7{,}065 \\
    \midrule
    Total & 21{,}249 \\
    \bottomrule
  \end{tabular}
\end{table}

\subsection{Dataset statistics summary}
\label{sec:dataset:stats}

\Cref{tab:dataset_stats} consolidates high-level statistics for quick reference.
Additional columns (frame counts, hours of video, camera types) require metadata not yet verified.

\begin{table}[t]
  \centering
  \caption{Dataset statistics (verified fields only). Extended columns are placeholders pending metadata confirmation.}
  \label{tab:dataset_stats}
  \begin{tabular}{ll}
    \toprule
    Statistic & Value \\
    \midrule
    Athletes & \numathletes{} \\
    Videos (attempts) & \numvideos{} \\
    Annotated frames & 35{,}825 \\
    Training phases (excl.\ unlabeled) & \numclasses{} \\
    Window size / stride & \windowsize{} / 1 \\
    Processed windows & \numwindows{} \\
    Feature dimension per frame & \featuredim{} \\
    Rebuild hash match (\texttt{X.npy}, \texttt{y.npy}) & Verified \\
    Native FPS / resolution & \textit{[TODO: metadata required]} \\
    Public release status & \textit{[Pending legal review]} \\
    \bottomrule
  \end{tabular}
\end{table}
